%%%%%%%%%%%%%%%%%%%%%%%%%%%%%%%%%%%%%%%%%%%%%%%%%%%%%%%%%%%%%%%%%%%%%%
%%%%%%%%%%%%%%%% Set up %%%%%%%%%%%%%%%%%%%%%%%%%%%%%%%%%%%%%%%%%%%%%%
%%%%%%%%%%%%%%%%%%%%%%%%%%%%%%%%%%%%%%%%%%%%%%%%%%%%%%%%%%%%%%%%%%%%%%
\usepackage[utf8]{inputenc}
\usepackage[english]{babel} % Language
\usepackage[T1]{fontenc}


%%%%%%%%%%%%%%%%%%%%%%%%% References and links %%%%%%%%%%%%%%%%%%%%%%%%%%%%%%%%%

\usepackage[hyperfootnotes=false]{hyperref} % Allows to hyperref references, footnotes, etc.
\hypersetup{colorlinks=true,citecolor={BrickRed}, linkcolor={BrickRed}, urlcolor={BrickRed}}  
%\usepackage[natbibapa]{apacite} % reference style
\usepackage[dvipsnames]{xcolor}
    \usepackage[authoryear]{natbib}
%%%%%%%%%%%%%%%%%%%%%%%%% Tables and figures %%%%%%%%%%%%%%%%%%%%%%%%%%%%%%%%%

\usepackage[flushleft]{threeparttable} % must-have for making table notes.
\usepackage{float} % used for putting pictures exactly where you want them. Use [H]
\usepackage{multirow,booktabs}
\usepackage{makecell}
\usepackage{graphics}
\usepackage{graphicx}
\usepackage{adjustbox} % Allows to adjust tables and figures.
\graphicspath{ {./images/} }
\usepackage{booktabs,caption}
\usepackage{comment}

%%%%%%%%%%%%%%%%%%%%%%%%% Math and symbols %%%%%%%%%%%%%%%%%%%%%%%%%%%%%%%%%

\usepackage{amsmath} % For math symbols
\usepackage{array} % For making arrays in math environments
\usepackage{amsthm} % For math symbols
\usepackage{amssymb} % For math symbols
\renewcommand\qedsymbol{Q.E.D.} % Changes the conclusion of a theorem. 
\theoremstyle{definition}
\newtheorem{definition}{Definition}[]
\newtheorem{theorem}{Proposition}
\usepackage{tocloft} % To change space in TOC
\usepackage{tikz} %allows for additional symbols such as \checkmark
\usepackage{caption}
\usepackage{subcaption}
\usetikzlibrary{shapes, arrows, positioning}

%%%%%%%%%%%%%%%%%%%%%%%%% Format and font %%%%%%%%%%%%%%%%%%%%%%%%%%%%%%%%%

\usepackage{geometry} % Allows to adjusts margins
\geometry{top = 25mm, bottom = 25mm, right = 25mm, left = 25mm}
\usepackage{titlesec} % Allows to customize contents in table of contents. 
\renewcommand{\abstractname}{\vspace{-\baselineskip}}
\renewcommand{\theequation}{\arabic{equation}} %\renewcommand{\theequation}{\arabic{section}.\arabic{equation}} if you want, for example Equation 2.1
%\usepackage{lmodern} %Use an usepackage to access fontfaces. 
%\fontfamily{lmr}\selectfont % Type the fontface
\usepackage{setspace} \setstretch{1.5} 
\usepackage{parskip} %this command will skip indent and make space between paragraphs.
\usepackage{pdfpages} % Allows for inserting pdf:s.
\usepackage{pdflscape} % Allows for rotating pages.
\usepackage[bottom]{footmisc} %sticks footnotes to the bottom of page
\interfootnotelinepenalty=10000 % makes footnotes stay on the same page and not continue on next.
\usepackage{csquotes}% Recommended for quotes

\renewcommand{\thesection}{\Roman{section}}
% \renewcommand{\thesubsection}{\Alph{subsection}} If you want only the subsection to be "A."
\renewcommand{\thesubsection}{\thesection.\Alph{subsection}}
\renewcommand{\thesubsubsection}{\thesubsection.\roman{subsubsection}}

\renewcommand{\thetable}{\Roman{table}}
\renewcommand{\thefigure}{\Roman{figure}}
\renewcommand{\theequation}{\arabic{equation}}
%\renewcommand{\theequation}{\arabic{section}.\arabic{equation}} if you want for example Equation 2.1

\providecommand{\tightlist}{%
  \setlength{\itemsep}{0pt}\setlength{\parskip}{0pt}}


%%%%%%%%%%%%%%%%%%%%%%%%%%%%%%%%%%%%%%%%%%%%%%%%%%%%%%%%%%%%%%%%%%%%%%%%%%
%%%%%%%%%%%%%%%%% End set up %%%%%%%%%%%%%%%%%%%%%%%%%%%%%%%%%%%%%%%%%%%%%
%%%%%%%%%%%%%%%%%%%%%%%%%%%%%%%%%%%%%%%%%%%%%%%%%%%%%%%%%%%%%%%%%%%%%%%%%%


